\begin{problem}[Warmup]{(a)}
    We have finally discussed some real numbers that are \emph{not} constructible! Give some examples.
\end{problem}
\unsolved

\begin{problem}[Warmup]{(b)}
    During our discussion of trisecting the angle, we proved that if an angle $\theta$ is constructible, then $\cos\theta$ is constructible. Explain why the converse of this statement is true too.
\end{problem}
\unsolved

\begin{problem}[Warmup]{(c)}
    Lambert proved in the 18th century that $\pi$ is not rational; why is this \emph{not} sufficient to show that the circle cannot be squared?
\end{problem}
\unsolved

\begin{problem}[Warmup]{(d)}
    Explain why it is possible to double the square.
\end{problem}
\unsolved

\begin{problem}[Warmup]{(e)}
    Explain why it is possible to `octuple' the cube.
\end{problem}
\unsolved

\begin{problem}{1}
    \begin{enumerate}
        \item[(a)] The number $\sqrt{2} + \sqrt{3}$ is constructible. It is an element of
        $F(\sqrt{3})$ where $F = \mathbb{Q}(\sqrt{2})$. In fact, $\sqrt{2} + \sqrt{3}$ is
        the root of the polynomial $x - (\sqrt{2} + \sqrt{3})$ in $F(\sqrt{3})$. Find a
        polynomial in $\mathbb{Q}[x]$ for which $\sqrt{2} + \sqrt{3}$ is a root.
        \item[(b)] Find a polynomial in $\mathbb{Q}[x]$ for which $\sqrt{4 + \sqrt{7}}$ is
        a root.
    \end{enumerate}
\end{problem}
\unsolved

\begin{problem}{2}
    Show that $\sqrt[3]{2} + \sqrt{2}$ is a root of the polynomial
    \[x^3 - 3\sqrt{2}x^2 + 6x - (2 + 2\sqrt{2})\]
    in $\mathbb{Q}(\sqrt{2})[x]$. Then obtain a sixth-degree polynomial in $\mathbb{Q}[x]$
    that has this number as a root.
\end{problem}
\unsolved

\begin{problem}{3}
    While it is impossible to square the circle, it is possible to square the rectangle.
    That is, given a rectangle, it is possible to construct a square of the same area. Do
    this.

    \emph{Hint:} Consider the diagram in the book (a semicircle on a diameter of length
    $a + b$, with a perpendicular of length $c$ erected at the point dividing the diameter
    into pieces $a$ and $b$). Show that $c^2 = ab$. Thus, if one constructed a square with
    sides of length $c$, it would have the same area as the given rectangle with sides $a$
    and $b$.
\end{problem}
\unsolved

\begin{problem}{4}
    \begin{enumerate}
        \item[(a)] Show how to square an arbitrary triangle.
        \item[(b)] Consider a \textbf{polygon} in the plane (by this, we mean just a
        bounded figure with edges that are line segments). How could you use part a to
        square such a figure?
    \end{enumerate}
\end{problem}
\unsolved

\begin{problem}{5}
    Suppose that $p$ and $q$ are elements of $F$, a subfield of the field of real numbers.
    Let $F(\sqrt{k})$ be a proper quadratic field extension of the field $F$. Prove that if
    the equation
    \[x^3 + px + q = 0\]
    has a solution in $F(\sqrt{k})$, then it has a solution in $F$.
\end{problem}
\unsolved

\begin{problem}{6}
    Use the trig identity
    \[\cos 3\theta = 4\cos^3\theta - 3\cos\theta,\]
    and the previous exercise to prove that the angle $2\pi/9$ is not constructible.
\end{problem}
\unsolved

\begin{problem}{7}
    Suppose that $n$ is a positive integer, with no integer cube root. Let $F(\sqrt{k})$ be
    a quadratic field extension of a field $F$. Prove that if
    $\sqrt[3]{n} \in F(\sqrt{k})$, then $\sqrt[3]{n} \in F$.
\end{problem}
\unsolved

\begin{problem}{8}
    Let $n$ be a positive integer, with no integer cube root. Use the previous exercise to
    prove that we cannot construct a cube with volume equal to $n$ times that of a given
    cube. (Thus, we cannot triple or quadruple the cube with compass and straightedge.)
\end{problem}
\unsolved

\begin{problem}{9}
    Consider a parabola (by analytic geometry, we may consider a curve of the form
    $f(x) = ax^2 + bx + c$). Pick two points
    \[P_1(x_1, f(x_1)) \text{ and } P_2(x_2, f(x_2))\]
    on the parabola, and consider the region bounded by the parabola and the line segment
    between the two points. This is called a \textbf{segment of the parabola}.
    \begin{enumerate}
        \item[(a)] Prove (using calculus) that the area of the segment is equal to
        $\frac{4}{3}$ times the area of the triangle $P_1P_2P_3$, where $P_3$ is the point
        on the parabola with $x$ coordinate $(x_1 + x_2)/2$.
        \item[(b)] Explain why this means that we can square the segment of the parabola.
        This result was known to Archimedes (ca. 250 B.C.); he proved it using the geometry
        of the parabola, rather than calculus.
    \end{enumerate}
\end{problem}
\unsolved

\begin{problem}{10}
    In the proof of Lemma 39.1 we needed to show that $x^3 - 2$ has only one real root; we
    did this using algebra. Prove this result using calculus.
\end{problem}
\unsolved

\begin{problem}{11}
    Suppose that $\theta$ is any fixed angle with positive radian measure. In this problem
    you will show that it is possible to construct an angle arbitrarily close in size to
    $\theta$.
    \begin{enumerate}
        \item[(a)] Suppose that $\epsilon$ is an arbitrary positive number. Describe how to
        construct an angle $\psi$ whose angle measure is smaller than $\epsilon$.
        \item[(b)] Explain why an integer multiple $n\psi$ of the angle constructed in part
        a must be larger than $\theta$.
        \item[(c)] Consider the \emph{smallest} positive integer $n$ so that $n\psi$ is
        larger than $\theta$ (why must $n$ exist?). Use $n$ to find an angle within
        $\epsilon$ radians of $\theta$.
        \item[(d)] Why does part c mean that we can come arbitrarily close to constructing
        the trisection of any given angle?
        \item[(e)] What is the philosophical difference between the construction in part d
        and the sought-for (impossible) trisection construction?
    \end{enumerate}
\end{problem}
\unsolved
